\section{Điều khung chưa giải quyết}
\label{sec:ch06-dieu-chua-giai-quyet}

\index{độ trung thực}%
\index{ground truth}%
Các trường hợp thu hồi được cho thấy khung đo được những giả định đã có, nhưng
không tự chọn giả định đúng cho một miền ứng dụng. Bài báo thử một cách chọn
bằng tối ưu, trên lớp mô hình mà ground truth có sẵn: mô hình tuyến tính
$f(y)=w^{\top}y$, ở đó tập đặc trưng thật là những $j$ có $|w_j|$ vượt một
ngưỡng. Attribution ở đây là toàn cục, không phụ thuộc $x$, và mỗi
\tn{độ đo}{measure} $\mu_j$ là một độ đo xác suất trên $[0,1]^d$. Với mỗi bộ
độ đo như thế, recall là tỉ lệ đặc trưng thật được attribution xếp trên một
ngưỡng thứ hai, và bài toán là tìm bộ độ đo làm recall lớn
nhất~\autocite{p17firstprinciples}. Chính điều kiện biết ground truth là giới
hạn đầu tiên mà bài báo tự nêu: không phải lúc nào ta cũng có tập đặc trưng
chuẩn để tối ưu theo cách đó.

Giới hạn thứ hai bài báo nêu là kỹ thuật, và nó nằm ở cách bài toán được giải.
Tối ưu trực tiếp trên không gian các độ đo nói chung khó và không có dạng
đóng, nên bài báo thay mỗi độ đo bằng trọng tâm của nó, điểm trung bình của
$y$ theo độ đo ấy như mục~\ref{sec:ch06-thu-hoi} đã dùng, và phát biểu điều
kiện tối ưu trên các trọng tâm: với mỗi đặc trưng thật $j$, trọng tâm $m$ của
$\mu_j$ phải thỏa $w^{\top}m\ge\alpha$ hoặc $w^{\top}m\le-\alpha$, với
$\alpha$ là ngưỡng xếp hạng, còn với các đặc trưng khác thì trọng tâm tự
do~\autocite{p17firstprinciples}. Phép thay độ đo bằng trọng tâm không song
ánh, vì nhiều độ đo khác nhau có cùng một trọng tâm. Vì vậy kết quả chỉ nói
trọng tâm tối ưu nằm ở đâu, chứ không mô tả hết các độ đo tối ưu, và cũng
không có cách đơn giản để đi ngược từ một trọng tâm về mọi độ đo sinh ra
nó~\autocite{p17firstprinciples}.

Còn một giới hạn thứ ba mà cuốn sách này thêm vào, nằm ở ý nghĩa của
attribution. Một độ đo có thể cho một phép tính chặt chẽ và một xấp xỉ chính
xác, nhưng vẫn đặt trọng lượng vào các đầu vào không phù hợp với câu hỏi của
người dùng hay với phân phối triển khai.
Vì vậy, khung này củng cố kết luận cuối chương~\ref{ch:attribution-theo-gradient}:
bảo đảm toán học của phép phân bổ chưa là bằng chứng cho độ trung thực của lời
giải thích.

Chương~\ref{ch:khung-hop-nhat} sẽ mở rộng câu hỏi từ attribution đặc trưng sang
một bản đồ có cả \tn{attribution dữ liệu}{data attribution} và
\tn{diễn giải cơ chế}{mechanistic interpretability}. Sau đó, chương~\ref{ch:do-do-trung-thuc} chuyển trọng tâm sang chính các dụng cụ dùng để đánh
giá lời giải thích. Khi ấy, một chỉ số cũng phải khai báo đối chiếu và giả định
của nó, giống như một độ đo trong chương này.

\index{độ đo|)}%
